\DocumentMetadata{
  pdfversion=2.0,
  pdfstandard=ua-2,
  lang=en-US,
  tagging=on,
  tagging-setup={math/setup=mathml-SE,tabsorder=structure}
}
\documentclass[11pt,letterpaper]{article}
\usepackage[margin=0.68in,includefoot]{geometry}
\usepackage{newcomputermodern}
\usepackage{amsmath,amscd,mathtools}
\usepackage{tikz,adjustbox}
\providecommand{\square}{\mdlgwhtsquare}
\usepackage[hidelinks]{hyperref}
\hypersetup{
  pdftitle={Numerical Linear Algebra First-Year Exam, August 2017},
  pdfauthor={University of Florida Department of Mathematics},
  pdfsubject={Graduate mathematics examination},
  pdfdisplaydoctitle=true,
  bookmarksopen=true
}
\setcounter{secnumdepth}{0}
\setlength{\parindent}{0pt}
\setlength{\parskip}{0.28em}
\setlength{\emergencystretch}{3em}
\setlength{\leftmargini}{2.15em}
\setlength{\leftmarginii}{2.65em}
\setlength{\leftmarginiii}{3em}
\renewcommand{\labelenumi}{\textbf{\theenumi.}}
\pagestyle{empty}
\raggedbottom
\allowdisplaybreaks
\newenvironment{examproblems}[1]{%
  \begin{enumerate}%
  \setcounter{enumi}{#1}%
  \setlength{\topsep}{0.35em}%
  \setlength{\partopsep}{0pt}%
  \setlength{\itemsep}{0.48em}%
  \setlength{\parsep}{0.18em}%
}{\end{enumerate}}
\newenvironment{examsubparts}{%
  \begin{enumerate}%
  \setlength{\topsep}{0.18em}%
  \setlength{\partopsep}{0pt}%
  \setlength{\itemsep}{0.16em}%
  \setlength{\parsep}{0.1em}%
}{\end{enumerate}}
\newenvironment{examinstructions}{%
  \par\begingroup\itshape
}{%
  \par\endgroup\vspace{0.18em}
}
\makeatletter
\renewcommand\section{\@startsection{section}{1}{0pt}%
  {-1.25ex plus -.25ex minus -.1ex}%
  {0.55ex plus .1ex}%
  {\normalfont\large\bfseries}}
\renewcommand{\@maketitle}{%
  \newpage\null\vskip -1em%
  {\footnotesize\bfseries
    \href{https://www.ufl.edu/}{University of Florida}, \href{https://math.ufl.edu/}{Department of Mathematics}\par}%
  \vskip -0.5em%
  \tagstructbegin{tag=Title}%
    {\LARGE\bfseries\href{https://gma.math.ufl.edu/exams/numerical-linear-algebra-first-year/}{Numerical Linear Algebra First-Year Exam}, August 2017\par}%
  \tagstructend%
  \vskip 0.25em%
  \tagmcbegin{artifact}%
    \hrule height 1pt%
  \tagmcend%
  \vskip 1em%
}
\makeatother
\title{Numerical Linear Algebra First-Year Exam, August 2017}
\author{}
\date{}
\begin{document}
\maketitle
\thispagestyle{empty}
\begin{examinstructions}
Do 4 (four) problems.
\end{examinstructions}
\begin{examproblems}{0}
\item
Assume that~\(A\) is Hermitian and all its eigenvalues are distinct and nonzero.
\begin{examsubparts}
\item[\textnormal{(a)}]
Show that each pair of distinct eigenvectors of~\(A\) are orthogonal.
\item[\textnormal{(b)}]
If~\(\lambda\) is an eigenvalue of~\(A\) with eigenvector~\(x\) with \(\lVert x\rVert_2=1\), then \(B=A-\lambda xx^*\) has the same eigenvectors as~\(A\) while~\(B\)'s eigenvalues are the same as those of~\(A\) except~\(\lambda\) is replaced with zero.
\end{examsubparts}
\item
This problem has two parts.
\begin{examsubparts}
\item[\textnormal{(a)}]
If~\(P\) is a projector, prove that \(\operatorname{null}(P)\cap\operatorname{range}(P)=\{0\}\) and \(\operatorname{null}(P)=\operatorname{range}(I-P)\).
\item[\textnormal{(b)}]
Prove that~\(P\) is an orthogonal projector if and only if it is Hermitian.
\end{examsubparts}
\item
This problem has three parts.
\begin{examsubparts}
\item[\textnormal{(a)}]
If both~\(A\) and~\(U\) are in \(\mathbb{C}^{m,m}\) and~\(U\) is unitary, prove that \(\lVert UA\rVert_F=\lVert AU\rVert_F=\lVert A\rVert_F\).
\item[\textnormal{(b)}]
If both~\(A\) and~\(U\) are in \(\mathbb{C}^{m,m}\) and~\(U\) is unitary, prove that \(\lVert UA\rVert_2=\lVert AU\rVert_2=\lVert A\rVert_2\).
\item[\textnormal{(c)}]
Prove that \(\lVert A\rVert_2=(\rho(A^*A))^{1/2}=\sigma_1\), where \(\sigma_1\) is the largest singular value of~\(A\).
\end{examsubparts}
\item
This problem has two parts.
\begin{examsubparts}
\item[\textnormal{(a)}]
If \(A\in\mathbb{C}^{m,n}\) with \(m\geq n\), prove that \(A^*A\) is invertible if and only if \(\operatorname{rank}(A)=n\).
\item[\textnormal{(b)}]
Give an explicit formula for \(\det(\lambda I-ww^*)\) when \(\lambda\in\mathbb{C}\), \(I\) is the \(m\times m\) identity matrix and \(w\in\mathbb{C}^m\).
\end{examsubparts}
\item
Assume that~\(T\) is tridiagonal and symmetric with the diagonal entries given by \(a_j\) for \(j=1,\ldots,m\) and the super- and sub-diagonal entries by \(b_j\) for \(j=1,\ldots,m-1\). Let \(p_k\) be the characteristic polynomial of the \(k\times k\) matrix in the upper left hand corner of~\(T\). Prove that 
\[
p_k(x)=(a_k-x)p_{k-1}(x)-b_{k-1}^2p_{k-2}(x).
\]
\end{examproblems}
\end{document}
